In Section~(\ref{golres}) we presented snapshots illustrating the realization of logic gates within a computer simulation of the Game of Life.
From this, one may conclude that a Boolean basis is effectively implemented within the system.
Once the principles for composing these gates into larger circuits are understood, increasingly complex computational structures can be constructed.

The existence of logic gates in the Game of Life implies that they form a functionally complete basis for Boolean logic.
Consequently, any Boolean function can be expressed as a finite composition of such gates.
This property is fundamental for the construction of general computational devices.
In particular, it is possible to design patterns that behave as a
\href{https://en.wikipedia.org/wiki/Finite-state_machine}{finite-state machine} coupled to two counters.
Such a system is computationally equivalent to a \href{https://en.wikipedia.org/wiki/Universal_Turing_machine}{universal Turing machine}

It follows that the Game of Life possesses the same theoretical computational power as any computer with unbounded
memory and no time constraints; that is, it is \href{https://en.wikipedia.org/wiki/Turing_complete}{Turing complete}~\cite{turingcomp}. This demonstrates that
complex dynamics can emerge solely from prescribed initial conditions and a simple update rule, forming a discrete dynamical system.




